\begin{afterword}
\summaryhead

The essay argues that humans, unlike other animals, can understand how their own minds produce beliefs, and can use that understanding to correct them. It begins with the physical chain by which light from untied shoelaces becomes a belief that they are untied. A mouse sees, but cannot know it has a visual cortex, so it cannot correct for optical illusions. Humans can, and science is this same capacity made systematic: reasoning about which processes make the mind mirror the world.

The example is a man in an online chatroom who expected no nuclear war for a hundred years and, asked why, answered ``pure hope.'' The author argues that hope is not evidence: across many possible worlds, optimists are not more likely to live in the peaceful ones. What makes you happy tells you about yourself, not about the world. A mind that knows its biases can apply second-order corrections, and this makes it far more powerful.

\responsehead

The essay claims that understanding the flaws of your mind lets you correct them. It names the correction and never performs one. The conclusion promises a ``reflective correction'' without saying what it is, how large it should be, or how anyone would know it had worked. The closing claim, that self-knowledge makes the lens ``far more powerful,'' comes with no evidence, although the author had warned five months earlier that knowing about biases can make people reason worse.

The essay's one worked case is thin. It is a chat with an anonymous man, undated, as the author recalls it. The man gave a reason about the world for the next twenty years; only the stretch beyond that rested on hope, and the essay treats the whole forecast as hope. The question the author posted, ``If no, why not?'', demanded reasons only from people who expected peace. The author's survey applied scrutiny to one answer, which is the fault the essay goes on to describe. From two words he then infers the man's unhappiness and a causal story about his brain.

The machinery around the example adds authority without adding content. The premise the essay needs, that a mind can be understood and corrected as a machine can, is stated in one sentence and defended with irony. The parallel worlds and the Tegmark footnote lend the weight of contested physics to a point that ordinary probability makes without them: optimism does not predict peace. The principle at the centre, that what you would like to be true is not evidence, is correct and familiar under the name of wishful thinking.

The one idea of substance, that a small effect should shift belief by a small amount, is in a footnote. It is the whole answer to the question of how large a correction should be, and the footnote announces a calculation that it never performs.

In short: an essay about correcting the mind's flaws that names the correction and never shows one, resting on a single recalled chat whose question applied scrutiny to one side only.
\end{afterword}
